\section{Project Grants}\label{project-grants}

On the description of the funding, for each project I show both the
global budget and the budget for the participating institution, in the
format (global budget / local budget).

\subsection{Project Grants as Global
Coordinator}\label{project-grants-as-global-coordinator}

\begin{enumerate}
\def\labelenumi{\arabic{enumi}.}
\item
  \ul{HAVATAR}
\end{enumerate}

Title: Healthcare Added-Value Applications of Tele-Autonomous Robots

Dates: 2021-03-01 to 2024-02-28 (successfully concluded)

Reference: \url{https://doi.org/10.54499/PTDC/EEI-ROB/1155/2020}

Funding: FCT (248K€ / 130K€)

\begin{quote}
Institutions: \emph{Instituto de Sistemas e Robótica, IST-ID}

\emph{Instituto de Sistemas e Robótica -- Universidade de Coimbra}

\emph{Hospital da Luz GLSMED Learning Health, S.A.}
\end{quote}

Role: Global Coordinator (Principal Investigator)

\begin{quote}
Abstract: The HAVATAR project developed new software skills for
telepresence robots that increase their level of autonomy, the
teleoperation comfort of the user, and its social presence at the remote
site. Two academic partners with long-standing international experience
in robotics projects (IST-ID Lisboa and ISR-UC Coimbra) will collaborate
for the development of the tele-autonomy skills and their integration in
the healthcare added-value applications, that will be tested in
realistic scenarios provided by Hospital da Luz Learning Health (HL.LH).
One application will be hosted at the HL.LH Simulation Center, where
practitioner students, health professionals or other actors will
remotely access a telepresence robot to dynamically observe simulated
clinical practices and interact with the local peers (e.g. simulated
surgical acts or simulated Intensive Care Unit in a pandemic scenario).
In this scenario, we will also study the issue of having multiple remote
users sharing the same local space and interacting with each other via
their proxy robots. The other application will deploy one telepresence
robot in the home of a patient and let the relatives or health
professionals interact and program the robot to acquire added value
variables and patient reported data. In this application, we will study
how extended periods of autonomy can be useful for healthcare
professionals.
\end{quote}

\begin{enumerate}
\def\labelenumi{\arabic{enumi}.}
\setcounter{enumi}{1}
\item
  \ul{FIREFRONT}
\end{enumerate}

\begin{quote}
Title: Real­-Time Forest Fire Mapping and Spread Forecast Using Unmanned
Aerial Vehicles
\end{quote}

Dates: 2019-03-01 to 2023-02-28 (successfully concluded)

Funding: FCT with reference PCIF/SSI/0096/2017 (377K€ / 102K€)

\begin{quote}
Institutions: \emph{Instituto de Sistemas e Robótica, IST-ID}

\emph{Aeroclube de Torres Vedras - Centro de Ciência do Ar e do Espaço}

\emph{Associação para o Desenvolvimento da Aerodinâmica Industrial}

\emph{Força Aérea Portuguesa}

\emph{Instituto de Telecomunicações}

\emph{UAVision, Engenharia de Sistemas, Ltd}
\end{quote}

Role: Global Coordinator (Principal Investigator)

\begin{quote}
Abstract: This project develops solution to support firefighting actions
in forest fires through the real­time detection and tracking of fire
fronts and re­burns. This is achieved by processing the information
acquired from manned and unmanned aerial vehicles equipped with
specialized sensing and communication systems, that overfly the affected
areas. This information can be made available to the coordination and
fighting units on a graphical interface the localization of the fire
events in georeferenced coordinates. Forecasts of the fire front
evolution, images of the fire area, wind speed and direction, and other
meteorological data of interest can also be supplied. This information
is of great value for decision making on firefighting actions.
\end{quote}

\begin{enumerate}
\def\labelenumi{\arabic{enumi}.}
\setcounter{enumi}{2}
\item
  \ul{VOAMAIS}
\end{enumerate}

\begin{quote}
Title: Computer Vision for the Operation of Unmanned Aerial Vehicles in
Maritime and WildfIre Scenarios
\end{quote}

Dates: 2019-02-01 to 2023-01-31 (successfully concluded)

Funding: FCT with reference PTDC/EEI-AUT/31172/2017 (222K€ / 154K€)

\begin{quote}
Institutions: \emph{Instituto de Sistemas e Robótica, IST-ID}

\emph{Força Aérea Portuguesa}

\emph{Marinha Portuguesa}
\end{quote}

Role: Global Coordinator (Principal Investigator)

\begin{quote}
Abstract: This project developed novel methods for target detection and
tracking in airborne and seaborne image sequences. The availability of
low-cost visual sensors (visible and IR) and the recent developments on
convolutional neural-networks and correlation filters are cost-effective
and energy-efficient solutions for detection and tracking of targets
capable of running onboard on air or sea vehicles. The results of the
project have a direct impact on strategic national activities such as
ocean and forest monitoring. To demonstrate the applicability of the
methods, three case scenarios are considered: (i) detection of forest
fires from airborne images, (ii) detection and tracking of sea vessels
from airborne and seaborne images, and (iii) detection, tracking and
pose estimation of aircrafts from seaborne images for teleguidance of
unmanned aerial vehicles (UAV's). These demonstrations will inform the
responsible institutions for the protection of the environment and
coastal operations of a novel technology to support their activities.
\end{quote}

\begin{enumerate}
\def\labelenumi{\arabic{enumi}.}
\setcounter{enumi}{3}
\item
  \ul{AHA}
\end{enumerate}

\begin{quote}
Title: Augmented Human Assistance
\end{quote}

Dates: 2014-08-01 to 2018-12-31 (successfully concluded)

Funding: FCT with reference CMUP-ERI/HCI/0046/2013 (590K€/307K€)

\begin{quote}
Institutions: \emph{Instituto de Sistemas e Robótica, IST-ID}

\emph{Faculdade de Motricidade Humana}, \emph{Universidade de Lisboa}

\emph{Madeira Interactive Technologies Institute M-ITI}

\emph{NOVA.ID.FCT}

\emph{PLUX, Engenharia de Biosensores, Lda}.

\emph{YDreams - Informática, SA}

\emph{Carnegie Mellon University} (US)
\end{quote}

Role: Global Coordinator (Principal Investigator)

\begin{quote}
Abstract: This project is a novel, integrative and cross-disciplinar
approach of 4 portuguese universities, CMU and 2 portuguese industry
partners that combines innovation and fundamental research in the areas
of human computer interaction, robotics, serious games and physiological
computing. AHA has developed a new generation of ICT based solutions
that have the potential to transform healthcare by optimizing resource
allocation, reducing costs, improving diagnoses and enabling novel
therapies, thus increasing quality of life. The project proposes the
development and deployment of a novel Robotic Assistance Platform
designed to support healthy lifestyle, sustain active aging, and support
those with motor deficits.
\end{quote}

\begin{enumerate}
\def\labelenumi{\arabic{enumi}.}
\setcounter{enumi}{4}
\item
  \ul{LIMOMAN}
\end{enumerate}

\begin{quote}
Title: Developmental Learning of Internal Models for Robotic
Manipulation based on Motor Primitives and Multisensory Integration
\end{quote}

Dates: 2014-05-01 to 2016-04-30 (successfully concluded)

Funding: EU-FP7-PEOPLE-2013-IEF, grant agreement 628315 (147K€)

\begin{quote}
Institutions: \emph{Instituto de Sistemas e Robótica, Instituto Superior
Técnico}
\end{quote}

Role: Global Coordinator

\begin{quote}
Abstract: LIMOMAN (developmental Learning of Internal MOdels for robotic
MANipulation based on motor primitives and multisensory integration)
addresses the key problem of improving the ability of current robots in
dexterous manipulation, inspired by three main aspects of the human
motor control system: internal models, learning and multisensory
integration. Internal models that encode the robot sensorimotor
capabilities and the dynamics of its interaction with the environment
are exploited to generate predictions that improve both robot perception
and motion control. These models are either acquired from scratch and/or
progressively refined and adapted through learning and optimization
techniques, using sensorimotor data collected by the robot during
goal-directed actions. The combination of different sensory modalities
(vision, proprioception, touch) and the use of probabilistic (Bayesian)
techniques permits to achieve robot behaviors that are effective, robust
and safe.
\end{quote}

\begin{enumerate}
\def\labelenumi{\arabic{enumi}.}
\setcounter{enumi}{5}
\item
  \ul{BIOMORPH}
\end{enumerate}

\begin{quote}
Title: Bio-inspired self-organizing robot computational morphologies
\end{quote}

Dates: 2014-02-01 to 2015-07-31 (successfully concluded)

Funding: FCT with reference EXPL/EEI-AUT/2175/2013 (49K€)

\begin{quote}
Institutions: \emph{Instituto de Sistemas e Robótica, IST-ID}
\end{quote}

Role: Coordinator (Principal Investigator)

\begin{quote}
Abstract: This project progressed towards the development of highly
efficient, scalable and adaptable control structures for autonomous
robotic agents, using biological principles for the creation of sparse
sensorimotor couplings and prediction mechanisms. The developed methods
were implemented and tested in concrete robotic test beds to prove the
proposed concepts in realistic scenarios and compared them with more
classical control approaches.
\end{quote}

\begin{enumerate}
\def\labelenumi{\arabic{enumi}.}
\setcounter{enumi}{6}
\item
  \ul{BIOLOOK}
\end{enumerate}

\begin{quote}
Title: Biomimetic Oculomotor Control for Humanoid Robots
\end{quote}

Dates: 2007-10-01 to 2010-11-30 (successfully concluded)

Funding: FCT with reference FCT - PTDC/EEA-ACR/71032/2006 (90K€)

\begin{quote}
Institutions: \emph{Instituto de Sistemas e Robótica, IST-ID}

\emph{University of Uppsala} (SW)
\end{quote}

Role: Global Coordinator (Principal Investigator)

\begin{quote}
Abstract: In this project we addressed the following scientific topics
in oculomotor control: (i) how to perform human-like eye-head movements
and postures in the execution of tasks and interaction with humans; (ii)
how oculomotor processes are learned and developed through childhood,
and how can be exploited to create adaptive autonomous systems; and
(iii) how robot behaviour can be modulated to implicitly communicate
emotions, intentions and goals to human users. The studied methodologies
were implemented and tested both in realistic dynamical simulations and
in prototype humanoid robotic platforms available in the consortium.
\end{quote}

\begin{enumerate}
\def\labelenumi{\arabic{enumi}.}
\setcounter{enumi}{7}
\item
  \ul{FOVEAR}
\end{enumerate}

\begin{quote}
Title: Foveal Vision for Robotics Applications
\end{quote}

Dates: 2006-04-01 to 2008-03-31 (successfully concluded)

Funding: FCT (Luso-Spain Integrated Action) with reference E-47/06 (4K€)

\begin{quote}
Institutions: \emph{Instituto de Sistemas e Robótica, Instituto Superior
Técnico}

\emph{Universidad Jaume I, Castellon de la Plana}
\end{quote}

Role: Principal Investigator

\begin{quote}
Abstract: The FOVEAR Project developed computer vision technology for
robotics applications with foveal vision sensors. These sensors yield a
significant reduction of the amount of information to process at any
time, while still allowing wide fields of view and high acuity in the
center of the visual field, like in the human retina. This allows the
real-time operation of robotic systems in generic and unconstrained
environments.
\end{quote}

\subsection{Project Grants As Coordinator of Participating
Institution}\label{project-grants-as-coordinator-of-participating-institution}

\begin{enumerate}
\def\labelenumi{\arabic{enumi}.}
\item
  \ul{ORIENT}
\end{enumerate}

\begin{quote}
Title: Goal-directed eye-head coordination in dynamic multisensory
environment.
\end{quote}

Dates: 2019-07-01 to 2022-12-31 (successfully concluded)

Reference: \url{https://doi.org/10.3030/693400}

Funding: H2020-EU.1.1, ERC-215-AdG, n. 693400 (2523K€/314K€)

\begin{quote}
Institutions: \emph{Stichting Radboud Universiteit, Nijmegen} (NL)

\emph{Instituto de Sistemas e Robótica, Instituto Superior Técnico}

\emph{Instituto de Sistemas e Robótica, IST-ID}
\end{quote}

Role: Participating Institution Coordinator

\begin{quote}
Abstract: Rapid object identification is crucial for survival of all
organisms but poses daunting challenges if many stimuli compete for
attention, and multiple sensory and motor systems are involved in the
processing, programming and generating of an eye-head gaze-orienting
response to a selected goal. How do normal and sensory-impaired brains
decide which signals to integrate (``goal''), or suppress
(``distracter'')? Audiovisual (AV) integration only helps for spatially
and temporally aligned stimuli. However, sensory inputs differ markedly
in their reliability, reference frames, and processing delays, yielding
considerable spatial-temporal uncertainty to the brain. Vision and
audition utilize coordinates that misalign whenever eyes and head move.
Meanwhile, their sensory acuities vary across space and time in
essentially different ways. As a result, assessing AV alignment poses
major computational problems, which so far have only been studied for
the simplest stimulus-response conditions.
\end{quote}

\begin{enumerate}
\def\labelenumi{\arabic{enumi}.}
\setcounter{enumi}{1}
\item
  \ul{ELEVAR}
\end{enumerate}

\begin{quote}
Title: Study of Vertical Structures with Robotized Aircrafts
\end{quote}

Dates: 2017-04-01 to 2019-03-31 (successfully concluded)

Funding: PT2020 - LISBOA-01-0247-FEDER-17924 (101K€)

\begin{quote}
Institutions: \emph{Tekever ASDS Lda}

\emph{Instituto de Sistemas e Robótica, Instituto Superior Técnico}

\emph{Laboratório Nacional de Engenharia Civil}
\end{quote}

Role: Participating Institution Coordinator

\begin{quote}
Abstract: The inspections and assessments of the structural integrity of
large infrastructures (eg. dams, bridges and high voltage cables) are
inadequate because they are not performed in a systematic manner and
they greatly depend on the personal criteria of the inspection
engineers. Additionally, certain areas of these structures or difficult
to reach by conventional means and the locations where data can be
acquired are limited which in turn can compromise the quality of the
gathered information and make these inspections costly. Some of the
already existing solutions based on multicopters are completely
dependent on a certified operator and/or GPS, lacking the ability to
perform an autonomous survey without the support of an infrastructure
that supplies the platform with global positioning signals. The solution
proposed by ELEVAR consists of utilizing an autonomous unmanned aircraft
integrated with sensors and vision-based navigation algorithms to solve
the above-mentioned problems.
\end{quote}

\begin{enumerate}
\def\labelenumi{\arabic{enumi}.}
\setcounter{enumi}{2}
\item
  \ul{iPEPE}
\end{enumerate}

\begin{quote}
Title: Cognitive Rehabilitation with Portable Exergame Platforms for the
Elderly
\end{quote}

Dates: 2020-06-01 to 2023-05-31 (successfully concluded)

Funding: P2020 with reference LISBOA-06-4234-FSE-000079 (30K€ / 74K€)

\begin{quote}
Institutions: \emph{Instituto das Irmãs Hospitaleiras do Sagrado Coração
de Jesus}

\emph{Instituto de Sistemas e Robótica, Instituto Superior Técnico}
\end{quote}

Role: Participating Institution Coordinator

\begin{quote}
Abstract: This project developed and tested augmented reality exergames
for the cognitive stimulation of the elder people living in nursing
homes or care institutions. These games put together physical and
cognitive tasks in entertaining and motivating activities and have been
shown to have benefits at the physical, cognitive, and social
dimensions. In this project a special attention is given to patient with
dementia, to assess if the proposed approach can be applied to this
fraction of the population.
\end{quote}

\begin{enumerate}
\def\labelenumi{\arabic{enumi}.}
\setcounter{enumi}{3}
\item
  \ul{SEAGULL}
\end{enumerate}

\begin{quote}
Title: Intelligent Systems for the Maritime Situational Awareness using
Unmanned Aerial Vehicles.
\end{quote}

Dates: 2013-07-01 to 2015-06-30 (successfully concluded)

Funding: QREN - FCOMP-01-0202-FEDER-034063 (870K€/97K€)

\begin{quote}
Institutions: \emph{Critical Software, SA}

\emph{Instituto de Sistemas e Robótica, Instituto Superior Técnico}

\emph{Faculdade de Engenharia da Universidade do Porto}

\emph{Força Aérea Portuguesa}

\emph{Marinha Portuguesa}
\end{quote}

Role: Participating Institution Coordinator

\begin{quote}
Abstract: Portugal is a maritime nation having the largest exclusive
economic zone (EEZ) of Europe. Upon acceptance of the proposal submitted
to the United Nations (UN), the Portuguese EEZ will be extended to an
area larger than 2 million square kilometres. In addition, Portugal is
known as a specialist in maritime matters not only by historical reasons
but also due to its contemporary activities on the sea. It is crucial
for Portugal to maintain the safety of human life at sea, maintain the
safety of its maritime area, and protect the sea and coastline natural
environments. The SEAGULL project consists on the use of unmanned air
vehicles (UAV's) as a not only viable but also cheaper solution to help
with these difficult tasks on such a large area. The objective is to
develop an intelligent system that couples the UAV with sensors like
visual, infra-red and hyperspectral cameras along with GPS receivers and
other sensors. Its purpose is to fly autonomously over the sea while
scanning it to automatically detect, identify, report and follow targets
like boats, shipwrecks, lifeboats, boats with suspicious behaviour, and
also hydrocarbyls and other chemical pollutants.
\end{quote}

\begin{enumerate}
\def\labelenumi{\arabic{enumi}.}
\setcounter{enumi}{4}
\item
  \ul{DICO(RE)2S}
\end{enumerate}

\begin{quote}
Title: Discount Coupon Recommendation and Redemption System
\end{quote}

Dates: 2011-07-01 to 2016-06-30 (successfully concluded)

Funding: EU-FP7-SME-262451 (1428K€/216K€)

\begin{quote}
Institutions: \emph{Fundacio Barcelona Media} (ES)

\emph{Nova Ventus Consulting SL} (ES)

\emph{Nextage SRL} (IT)

\emph{Aitek SPA} (IT)

\emph{Observit -- Tecnologias de Visão por Computador, Lda}

\emph{Instituto de Sistemas e Robótica, Instituto Superior Técnico}
\end{quote}

Role: Participating Institution Coordinator

\begin{quote}
Abstract: DICO(RE)2S is an SME European Consortium devoted to develop
and deploy a coupon-based discount campaign platform that will exploit
state-of-the-art information and communication technologies and
processes for providing consumers and retailers/manufactures with a
personalized and efficient environment for maximizing customer
satisfaction and business profitability. The distributed network of
coupon redemption points will implement pattern recognition systems for
reading and validating coupons as well as security encryption and
transmission systems that will guarantee the integrity of the overall
coupon redemption strategy as well as the reliability of the data.
\end{quote}

\begin{enumerate}
\def\labelenumi{\arabic{enumi}.}
\setcounter{enumi}{5}
\item
  \ul{HDA}
\end{enumerate}

\begin{quote}
Title: High-Definition Analytics.
\end{quote}

Dates: 2010-05-01 to 2013-10-30 (successfully concluded)

Funding: QREN - FCOMP-01-0202-FEDER-013750 (536K€/203K€)

\begin{quote}
Institutions: \emph{Observit -- Tecnologias de Visão por Computador,
Lda}

\emph{Instituto de Sistemas e Robótica, Instituto Superior Técnico}
\end{quote}

Role: Participating Institution Coordinator and Global Scientific
Coordinator

\begin{quote}
Abstract: HD-Analytics is a QREN research project aimed at the
development of new algorithms and methods to address important video
analytics challenges, in the context of new technology available in the
IP camera market: high-resolution multi-megapixel cameras and
high-resolution omni-directional cameras. Thus, new algorithms and
methods are required to deal with the increased amount of data.
\end{quote}

\begin{enumerate}
\def\labelenumi{\arabic{enumi}.}
\setcounter{enumi}{6}
\item
  \ul{HANDLE}
\end{enumerate}

\begin{quote}
Title: Developmental pathway towards autonomy and dexterity in robot
in-hand manipulation
\end{quote}

Dates: 2009-02-01 to 2013-01-31 (successfully concluded)

Funding: EU-FP7-ICT-231640 (8190K€/749K€)

\begin{quote}
Institutions: \emph{Universite Pierre et Marie Curie, Paris 6} (FR)

\emph{Universitaet Hamburg} (DE)

\emph{Universidad Carlos III de Madrid} (ES)

\emph{Commissariat a L'Energie Atomique et Aux Energies Alteratives}
(FR)

\emph{Instituto de Sistemas e Robótica, Instituto Superior Técnico}

\emph{Faculdade de Ciências e Tecnologia da Universidade de Coimbra}

\emph{Orebro University} (SW)

\emph{Kings College London} (UK)

\emph{The Shadow Robot Company Limited} (UK)
\end{quote}

Role: Participating Institution Coordinator and Global Technical Manager

\begin{quote}
Abstract: The HANDLE project aims at understanding how humans perform
the manipulation of objects to replicate grasping and skilled in-hand
movements with an anthropomorphic artificial hand, and thereby move
robot grippers from current best practice towards more autonomous,
natural, and effective articulated hands. The project implies not only
focusing on technological developments but also working with fundamental
multidisciplinary research aspects to endow the robotic hand system with
advanced perception capabilities, high level feedback control and
elements of intelligence that allow recognition of objects and context,
reasoning about actions and a high degree of recovery from failure
during the execution of dexterous tasks.
\end{quote}

\begin{enumerate}
\def\labelenumi{\arabic{enumi}.}
\setcounter{enumi}{7}
\item
  \ul{VEMUCARV}
\end{enumerate}

\begin{quote}
Title: Spatial Validation of Complex Urban Grids in Virtual Imersive
Environments
\end{quote}

Dates: 2009-02-01 to 2013-01-31 (successfully concluded)

Funding: FCT - POCTI/AUR/48123/2002 (25K€)

\begin{quote}
Institutions: \emph{Institute of Mechanical Engineering, Instituto
Superior Técnico}

\emph{Instituto de Sistemas e Robótica, Instituto Superior Técnico}

\emph{Câmara Municipal de Lisboa}
\end{quote}

Role: Participating Institution Coordinator

\begin{quote}
Abstract: The main goals of this project are related to the
semi-automatic acquisition and maintenance of 3D virtual reality models
of urban areas. It is intended to use registered aerial images and low
altitude laser range scans to acquire 3D data of city structure. This
data will be processed in order to segment relevant structures for urban
planning (buildings, roads, green areas, etc). Range information
provides a very rich description of 3D structure but lack photometric
information. Aerial photos provide this information, allowing to
pre-segment regions based on color and texture. The main scientific
innovation of this project is the combined use of 2D (aerial images) and
3D (range scans) to simplify and improve the building extraction
process. Most current approaches use one or the other types of data
exclusively. The final result provides a computer model which stands for
a mix geometry-image database that can interface to GIS software
available (at CML), as well as the generation of real-time walkthrough
with thematic information. The results of this project are to be
integrated on Lisbon City Hall public computational facilities.
\end{quote}

\subsection{Project Grants As Team
Member}\label{project-grants-as-team-member}

\begin{enumerate}
\def\labelenumi{\arabic{enumi}.}
\item
  \ul{RePAIR}
\end{enumerate}

\begin{quote}
Title: Reconstructing the Past: Artificial Intelligence and Robotics
Meet Cultural Heritage
\end{quote}

Dates: 2021-06-01 to 2024-11-30 (ongoing)

Funding: H2020 FETOPEN - Grant agreement 964854 (3522K€/575K€)

\begin{quote}
Institutions: \emph{Universita Ca'Foscari, Venezia} (IT)

\emph{Instituto de Sistemas e Robótica, IST-ID}

\emph{Ben-Gurion University of the Negev} (IR)

\emph{Fondazione Istituto Italiano di Tecnologia} (IT)

\emph{Rheinische Friedrich-Wilhelms-Universitat Bonn} (DE)

\emph{Ministero Della Coltura} (IT)
\end{quote}

Role: Team Member

\begin{quote}
Abstract: The physical reconstruction of shattered artworks is one of
the most labour-intensive steps in archaeological research. Dug out from
excavation sites are countless ancient artefacts, such as vases,
amphoras and frescoes, that are damaged. The EU-funded RePAIR project
will facilitate the reconstruction process to bring ancient artworks
back to life. Specifically, it will develop an intelligent robotic
system that can autonomously process, match and physically assemble
large fractured artefacts in a fraction of the time required by humans.
This new system will be tested on iconic case studies from the UNESCO
World Heritage Site of Pompeii. It will restore two world-renowned
frescoes, which are in thousands of broken pieces and currently in
storerooms.
\end{quote}

\begin{enumerate}
\def\labelenumi{\arabic{enumi}.}
\setcounter{enumi}{1}
\item
  \ul{ShiftHRI}
\end{enumerate}

\begin{quote}
Title: Exploring the Transfer of Agency to Older Adults in HRI
\end{quote}

Dates: 2022-03-01 to 2023-05-28 (successfully concluded)

Funding: FCT, CMU/TIC/0026/2021 (63K€/26K€)

\begin{quote}
Institutions: \emph{LASIGE}

\emph{Instituto de Sistemas e Robótica, IST-ID}

\emph{Carnegie Mellon University} (US)
\end{quote}

Role: Team Member

\begin{quote}
Abstract: Our vision is to advance research on human-robot interaction
to facilitate a new type of relationship between elders and assistive
robots. This relationship will be built on a deep understanding of human
values and what is critical to older adults. Our previous work shows
that values are manifested through interactions with products, and that
this becomes more critical for elders as they experience physical,
cognitive, and emotional decline. In addition, what elders value can
change dynamically, even day to day, as elders experience decline. In
ShiftHRI, we will start by exploring important values and co-designing
HRI scenarios that help people to fulfil them. This paradigm shift has
the potential to significantly impact the research done in the area.
Contrary to the state of the art, where most robots have strict
predefined user interfaces and routines, shifting agency to older adults
brings several challenges to the control of those robots and in the
service they offer.
\end{quote}

\begin{enumerate}
\def\labelenumi{\arabic{enumi}.}
\setcounter{enumi}{2}
\item
  \ul{IntelligentCare}
\end{enumerate}

\begin{quote}
Title: Intelligent Multimorbidity Management System
\end{quote}

Dates: 2020-04-01 to 2023-06-30 (successfully concluded)

Funding: LISBOA-01-0247- FEDER-045948 (3522K€/497K€)

\begin{quote}
Institutions: \emph{GLSMED Learning Health S.A.}

\emph{Instituto de Sistemas e Robótica, Instituto Superior Técnico}

\emph{Instituto de Sistemas e Robótica, IST-ID}

\emph{Hospital da Luz S.A.}

\emph{Instituto De Engenharia De Sistemas E Computadores, Investigação e
Desenvolvimento Em Lisboa}

\emph{Priberam Informática S.A.}

\emph{Carnegie Mellon University} (US)
\end{quote}

Role: Team Member

\begin{quote}
Abstract: This project focus on managing the issue of a growing
population with multimorbidity (co-occurring diseases), a worldwide
problem to the sustainability of the healthcare system. IntelligentCare
aims to develop a patient-centric solution to manage multimorbidity
conditions by using analytical methods to explore data from the
electronic health records and measures reported remotely by the
patients. Under the value-based healthcare framework, the project will
develop methods for early signalization and characterization of
patients, identification of personalized clinical pathways based on
evidence, and based on patient value vision, together with monitoring
and prediction of hospital interactions. To achieve these goals research
will be conducted in advanced analytical algorithms, particularly deep
learning methods, to fully explore data. This is an innovative approach
to evaluate the multimorbidity condition by introducing the concept of
value to the patient in the process of characterization and prediction
of patient interactions with the hospital.
\end{quote}

\begin{enumerate}
\def\labelenumi{\arabic{enumi}.}
\setcounter{enumi}{3}
\item
  \ul{ASTRIIS}
\end{enumerate}

\begin{quote}
Title: Atlantic Sustainability Through Remote and In-situ Integrated
Solutions

Dates: 2020-03-01 to 2023-06-30 (successfully concluded)

Funding: POCI-01-0247-FEDER-046092 (891K€/273K€)

Institutions: \emph{Tekever}

\emph{+Atlantic CoLAB}

\emph{CEiiA}

\emph{Hidromod, modelação em engenharia}

\emph{Institute for Systems and Robotics, IST}

\emph{ISQ}

\emph{LSTS, Underwater Systems and Technology Laboratory, FEUP}

\emph{MARETEC, Marine Environment and Technology Center, IST}

\emph{CERENA, Center of Natural Resources and Environment, IST}

\emph{Ocean Infinity}

\emph{Ocean Scan, Marine Systems \& Technology Lda}

\emph{Spin.Works, S.A.}

\emph{WavEC Offshore Renewables}

\emph{Universidade do Minho}

\emph{Universidade do Algarve}

Role: Team Member

Abstract: This project is a Portuguese effort to support the sustainable
growth of the blue economy by developing technical and scientific
knowledge for the design, demonstration and implementation of integrated
and customizable information products and services. The project focuses
on sectors of the blue economy with high potential for value creation,
such as marine renewable energy, aquaculture, coastal tourism, or
recreational boating, and is oriented to the needs of users, which
include investors, regulators, policymakers, consumers, and scientific
researchers. The ISR team participated in the development of remote
sensing solutions for environmental monitoring of the Atlantic national
coastal waters.
\end{quote}

\begin{enumerate}
\def\labelenumi{\arabic{enumi}.}
\setcounter{enumi}{4}
\item
  \ul{RBCog-Lab}
\end{enumerate}

\begin{quote}
Title: Robotics, Brain and Cognition Lab
\end{quote}

Dates: 2017-06-01 to 2021-11-30 (successfully concluded)

\begin{quote}
Funding: FCT - Portuguese Roadmap of Research Insfrastructures
01/SAICT/2016 SAICT Proj. 22084 (300K€)

Institutions: \emph{Instituto de Sistemas e Robótica, Instituto Superior
Técnico}

\emph{Instituto de Sistemas e Robótica, IST-ID}
\end{quote}

Role: Team Member

\begin{quote}
Abstract: The Robotics, Brain and Cognition Lab (RBCogLab) is part of
the Portuguese Roadmap of Research Infrastructures. RBCog-Lab supports
an integrated multidisciplinary research endeavour (neuroscience,
developmental psychology, cognitive science, psychophysics, machine
learning and other engineering subjects) not only promoting novel
robotic methodologies but also contributing to a better understanding of
the human brain. To be effective and generate a positive impact in the
society, this multidisciplinary research needs to be supported by
advanced equipment and specialized technical expertise, that might not
be available in small research and academic centres with low budgets.
The RBCogLab research infrastructure aims at filling this gap by
providing the scientific community with a centralized open laboratory
resource for robotics experiments and human measurements, offering both
state of the art devices and specialized know-how.
\end{quote}

\begin{enumerate}
\def\labelenumi{\arabic{enumi}.}
\setcounter{enumi}{5}
\item
  \ul{ALHTOUR}
\end{enumerate}

\begin{quote}
Title: Assisted Living Technologies for the Health Tourism Sector
\end{quote}

Dates: 2016-01-01 to 2018-12-31 (successfully concluded)

Funding: H2020-TWINN-2015 Grant Agreement 692311 (1175K€/40K€)

\begin{quote}
Institutions: \emph{Universidade de Lisboa}

\emph{Faculdade de Motricidade Humana}

\emph{Faculdade de Medicina da Universidade de Lisboa}

\emph{Faculdade de Ciências da Universidade de Lisboa}

\emph{Instituto de Ciências Sociais}

\emph{Instituto Superior Técnico}

\emph{Katholieke Universiteit Leuven} (BE)

\emph{Universiteit Maastricht} (NL)

\emph{Universita degli studi di Macerata} (IT)

Role: Coordinator of activities at the Institute for Systems and
Robotics, Instituto Superior Técnico.

Abstract: The ALHTOUR project aims to strengthen and stimulate
scientific excellence, and innovation capacity in technologies for
independent living, applied to the health tourism at UL (Portugal), by
preparing for the set-up of a ``Health Tourism Living Lab'', identified
as a key driver for territorial development.
\end{quote}

\begin{enumerate}
\def\labelenumi{\arabic{enumi}.}
\setcounter{enumi}{6}
\item
  \ul{SPARSIS}
\end{enumerate}

\begin{quote}
Title: Sparse Modeling and Estimation of Motion Fields
\end{quote}

Dates: 2016-07-01 to 2019-12-31 (successfully concluded)

Funding: FCT - PTDC/EEI-PRO/0426/2014 (199K€)

\begin{quote}
Institutions: \emph{Instituto de Sistemas e Robótica, Associação do
Instituto Superior Técnico para a Investigação e o Desenvolvimento}

\emph{Instituto De Engenharia De Sistemas E Computadores, Investigação e
Desenvolvimento Em Lisboa}

\emph{Instituto de Telecomunicações}
\end{quote}

Role: Workpackage Coordinator

\begin{quote}
Abstract: In this project we will explore the use of sparse techniques
to improve the estimation of multiple motion fields as well as space
varying switching matrix (stochastic matrix) that describes the
switching process associated to the movement of each target. As
scientific outcomes of the project, we expect to obtain reliable
estimates for the model parameters (many hundreds), to remove the over
smoothed character of the field estimates. We also expect to speed up
the estimation procedure bringing it closer to real time applications
and extend these techniques to multicamera settings. We also expect to
develop an adaptive algorithm for the online estimation of multiple
motion fields able to update in a recursive way the number of fields and
the field estimates whenever new information arrives.
\end{quote}

\begin{enumerate}
\def\labelenumi{\arabic{enumi}.}
\setcounter{enumi}{8}
\item
  \ul{INSIDE}
\end{enumerate}

\begin{quote}
Title: Intelligent Networked Robot Systems for Symbiotic Interaction
with Children with Impaired Development
\end{quote}

Dates: 2014-08-01 to 2018-07-31 (successfully concluded)

Funding: FCT with reference CMUP­ERI/HCI/0051/2013 (590K€/964K€)

\begin{quote}
Institutions: \emph{INESC-ID}

\emph{ISR-Lisboa, LARSyS, IST-ID}

\emph{Hospital Garcia de Horta EPE}

\emph{IDMIND -- Engenharia de Sistemas, LDA}

\emph{PLUX, Engenharia de Biosensores, LDA}

\emph{VoiceInteraction -- Tecnologias de Processamento da Fala, S.A.}

\emph{Carnegie Mellon University} (US)
\end{quote}

Role: Team Member

\begin{quote}
Abstract: The INSIDE initiative explores symbiotic interactions between
humans and robots in joint cooperative activities, and addresses the
following key research problems: (i) How can robots plan their course of
action to coordinate with and accommodate for the actions of their human
teammates? (ii) How can task, context and environment information
collected from a set of networked sensors (such as cameras and biometric
sensors) be exploited to create more natural and engaging interactions
between humans and robots involved in a joint cooperative activity in a
physical environment? INSIDE developed new hardware and software
solutions that support the aforementioned research towards the
deployment of a real-world interaction scenario where children with
impaired development interact with a networked robot system in a joint
cooperative task with therapeutical purposes.
\end{quote}

\begin{enumerate}
\def\labelenumi{\arabic{enumi}.}
\setcounter{enumi}{7}
\item
  \ul{Poeticon++}
\end{enumerate}

\begin{quote}
Title: Robots Need Language: A Computational Mechanism for
Generalisation and Generation of New Behaviours in Robots
\end{quote}

Dates: 2012-01-01 to 2015-12-31 (successfully concluded)

Funding: EU-FP7-ICT-288382 (3880K€/470K€)

\begin{quote}
Institutions: \emph{Fondazione Istituto Italiano do Tecnologia} (IT)

\emph{Instituto de Sistemas e Robótica, Instituto Superior Técnico}

\emph{Cognitive Systems Research Institute} (GR)

\emph{Athina} (GR)

\emph{University of Plymouth} (UK)

\emph{The University of Maryland Foundation, Inc} (US)
\end{quote}

Role: Deputy Scientific Coordinator at Participating Institution

\begin{quote}
Abstract: Handling novel situations beyond learned schemas or set
behaviours is still a quest in engineering cognitive and embodied
systems. Powerful generalisation mechanisms are necessary for any agent
to operate effectively in real-world environments. POETICON++ suggests
that natural language can be used as a learning tool for generalisation
of learned behaviours \& perceptual experiences and generation of new
behaviours \& experiences (creativity). The main objective of POETICON++
is the development of a computational mechanism for such generalisation
of motor programs and visual experiences for robots. To this end, it
will integrate natural language and visual action/object recognition
tools with motor skills and learning abilities.
\end{quote}

\begin{enumerate}
\def\labelenumi{\arabic{enumi}.}
\setcounter{enumi}{8}
\item
  \ul{AuReRo}
\end{enumerate}

\begin{quote}
Title: Human-robot interaction with field robots using augmented reality
and interactive mapping
\end{quote}

Dates: 2011-04-01 to 2014-09-30 (successfully concluded)

Funding: FCT - PTDC/EIA-CCO/113257/2009 (130K€)

\begin{quote}
Institutions: \emph{Instituto de Sistemas e Robótica, Instituto Superior
Técnico}

\emph{Madeira Interactive Technologies Institute M-ITI}
\end{quote}

Role: Co-Principal Investigator

\begin{quote}
Abstract: To achieve effective human-robot interaction (HRI), we propose
presenting a pair of stereo camera feeds from a robot to an operator
through an Augmented Reality (AR) head-mounted display (HMD). This will
provide an immersive experience to the controller and has the capacity
to display rich contextual information. To further enhance this display,
we will superimpose mapping data generated by SLAM-6D (Simultaneous
Localization And Mapping) algorithms over the video input using
Augmented Reality (AR) techniques.
\end{quote}

\begin{enumerate}
\def\labelenumi{\arabic{enumi}.}
\setcounter{enumi}{9}
\item
  \ul{FirstMM}
\end{enumerate}

\begin{quote}
Title: Flexible Skill Acquisition and Intuitive Robot Tasking for Mobile
Manipulation in the Real World
\end{quote}

Dates: 2010-02-01 to 2013-07-31 (successfully concluded)

Funding: EU-FP7-ICT-248258 (2997K€/351K€)

\begin{quote}
Institutions: \emph{Albert-Ludwigs-Univeritaet Freiburg,} (DE)

\emph{Instituto de Sistemas e Robótica, Instituto Superior Técnico}

\emph{Ecole Polytechnique Federale de Lausanne} (CH)

\emph{Kuka Laboratories GmbH} (DE)

\emph{Fraunhofer Institute} (DE)

\emph{Technische Universitat Berlin} (DE)

\emph{Idryma Technologias Kai Erevnas} (DE)
\end{quote}

Role: Deputy Scientific Coordinator at Participating Institution

\begin{quote}
Abstract: Many industrial processes highly depend on the reliability and
robustness of robotic manipulators. Research on mobile robots has led to
systems that demonstrated the capability of safe and accurate
navigation. The goal of the FIRST-MM project is to integrate these two
areas in the context of a real-world application scenario. The project
will build upon and extend recent results in robot programming,
navigation, manipulation, perception, learning by instruction, and
statistical relational learning to develop advanced technology for
autonomous mobile manipulation robots that can flexibly be instructed
even by non-expert users to perform challenging manipulation and
transportation tasks in real-world environments.
\end{quote}

\begin{enumerate}
\def\labelenumi{\arabic{enumi}.}
\setcounter{enumi}{10}
\item
  \ul{RoboSoM}
\end{enumerate}

\begin{quote}
Title: A Robotic Sense of Movement
\end{quote}

Dates: 2009-12-01 to 2013-05-31 (successfully concluded)

Funding: EU-FP7-ICT-248366 (1659K€/391K€)

\begin{quote}
Institutions: \emph{Scuola Superiore di Studi S'Anna,} (IT)

\emph{Instituto de Sistemas e Robótica, Instituto Superior Técnico}

\emph{Centre National de la Recherche Scientifique} (FR)

\emph{Waseda University} (JP)

\emph{College de France} (FR)
\end{quote}

Role: Deputy Scientific Coordinator at Participating Institution

\begin{quote}
Abstract: The objective of RoboSoM is to investigate and to apply new
approaches to the design and development of humanoid robots with
advanced perception and action capabilities, showing robust, adaptive,
predictive, and effective behaviour in the real world. The expected
robot behaviour is the capability to follow a visual target by
coordinating eye, head, and leg movements, with head stabilization,
walking smoothly and effectively in an unstructured environment, with a
robust reactive behaviour, improved by predictions.
\end{quote}

\begin{enumerate}
\def\labelenumi{\arabic{enumi}.}
\setcounter{enumi}{11}
\item
  \ul{MMCACC}
\end{enumerate}

\begin{quote}
Title: Modern Monte Carlo Algorithms for Computational Control
\end{quote}

Dates: 2007-09-01 to 2010-12-31 (successfully concluded)

Funding: FCT - PTDC/EEA-ACR/70174/2006 (80K€)

\begin{quote}
Institutions: \emph{Instituto de Sistemas e Robótica, Instituto Superior
Técnico}

\emph{University of British Columbia} (CA)
\end{quote}

Role: Team Member

\begin{quote}
Abstract: The objective is to develop modern Monte Carlo methods to
solve discrete time stochastic optimal control problems in both the
fully and partially observed cases for nonlinear and/or non-Gaussian
models. It aims to combine modern computational tools developed actively
in statistics (Markov chain Monte Carlo, Sequential Monte Carlo aka
Particle filters) to ideas developed in automatic control, operation
research (gradient estimation) and reinforcement learning (value
function and policy parameterization). The focus is in developing
generic algorithms that can be applied in different domains and provide
a unified framework for this type of problems.
\end{quote}

\begin{enumerate}
\def\labelenumi{\arabic{enumi}.}
\setcounter{enumi}{12}
\item
  \ul{URUS}
\end{enumerate}

\begin{quote}
Title: Ubiquitous Networking Robotics in Urban Settings
\end{quote}

Dates: 2006-12-01 to 2009-11-30 (successfully concluded)

Funding: EU-FP6-IST-045062 (2600K€/80K€)

\begin{quote}
Institutions: \emph{Universitat Politecnica de Catalunya} (ES)

\emph{Instituto de Sistemas e Robótica, Instituto Superior Técnico}

\emph{Agencia d'Ecologia Urbana de Barcelona} (ES)

\emph{Associacion de la Investigacion Y Cooperacion Industrial de
Andalucia ``F. de Paula Rojas''} (ES)

\emph{Centre National de la Recherche Scientifique} (FR)

\emph{ETH Zurich} (CH)

\emph{Robotech SRL} (IT)

\emph{Scuola Superiore di Studi S'Anna} (IT)

\emph{Telefonica Investigacion Y Desarrollo AS Unipersonal} (ES)

\emph{The University of Surrey} (UK)

\emph{Universidad de Zaragoza} (ES)
\end{quote}

Role: Scientific Coordinator of Computer Vision at Participating
Institution

\begin{quote}
Abstract: European ancient cities are becoming difficult places to live
due to noise, pollution, lack of quality facilities and security.
Moreover, the average age of people living in large European cities is
growing and in a short period of time there will be an important
community of elderly people. City Halls are becoming conscious of this
problem and are studying solutions, for example by reducing the areas of
free circulation of cars. Free car areas imply a revolution in the
planning of urban settings, for example, by imposing new means for
transportation of goods to the stores, security issues, new ways of
human assistance, etc. In this project we want to analyse and test the
idea of incorporating a network of robots (which include humans, robots,
intelligent sensors, intelligent devices and communications) in order to
improve quality of life in such urban areas.
\end{quote}

\begin{enumerate}
\def\labelenumi{\arabic{enumi}.}
\setcounter{enumi}{13}
\item
  \ul{CONTACT}
\end{enumerate}

\begin{quote}
Title: Learning and Development of Contextual Action
\end{quote}

Dates: 2005-09-01 to 2009-02-28 (successfully concluded)

Funding: EU-FP6-NEST-5010 (NA/380K€)

\begin{quote}
Institutions: \emph{Universita degli studi di Genova,} (IT)

\emph{Instituto de Sistemas e Robótica, Instituto Superior Técnico}

\emph{Universita degli studi di Ferrara} (IT)

\emph{Uppsala Universitet} (SW)

\emph{University of Stockholm} (SW)
\end{quote}

Role: Workpackage Scientific Coordinator at Participating Institution

\begin{quote}
Abstract: By spontaneously generating, observing and imitating
manipulative gestures, infants quickly learn how to interact with their
environment. At the same time, communicative abilities develop leading
to a more and more complex use of speech. Are these motor abilities
developing independently? Or are there fundamentally similar mechanisms
in the development of perception and production, for both speech and
manipulation? We believe that the later is the case. This project
represents an ambitious attempt to investigate this parallel development
of manipulatory and speech related gestures from a multi-disciplinary
perspective. Experts in Robotics, Neuroscience, and Child-Development
will tightly collaborate to build an artifact (an embodied system) which
will develop its motor capabilities both in manipulation and speech.
\end{quote}

\begin{enumerate}
\def\labelenumi{\arabic{enumi}.}
\setcounter{enumi}{14}
\item
  \ul{Robot-Cub}
\end{enumerate}

\begin{quote}
Title: Robotic Open-architecture Technology for Cognition, Understanding
and Behaviour
\end{quote}

Dates: 2004-09-01 to 2010-01-31 (successfully concluded)

Funding: EU-FP6-IST-2004-004370 (8500K€/594K€)

\begin{quote}
Institutions: \emph{Universita degli studi di Genova} (IT)

\emph{Universita degli studi di Ferrara} (IT)

\emph{Instituto de Sistemas e Robótica, Instituto Superior Técnico}

\emph{Uppsala Universitet} (SW)

\emph{Ecole Polytechnique Federale de Lausanne} (CH)

\emph{University of Zurich} (CH)

\emph{Fondazione European Brain Research Institute Rita Levi} (IT)

\emph{Scuola Superiore di Studi S'Anna} (IT)

\emph{Fondazione Istituto Italiano do Tecnologia} (IT)

\emph{Telerobot S.R.L.} (IT)

\emph{The University of Hertfordshire Higher Education Corporation} (UK)

\emph{The University of Salford} (UK)

\emph{The University of Sheffield} (UK)
\end{quote}

Role: Workpackage Scientific Coordinator at Participating Institution

\begin{quote}
Abstract: The main goals of RobotCub are (i) to create an open robotic
platform for embodied research that can be taken up and used by the
research community at large to further their approach to the development
of humanoid-based cognitive systems and (ii) to advance our
understanding of several key issues in cognition by exploiting this
platform in the investigation of cognitive capabilities. The scientific
objective of RobotCub is, therefore, to jointly design the mindware and
the hardware of a humanoid platform to be used to investigate human
cognition and human-machine interaction. We call this platform CUB or
Cognitive Universal Body. It is worth remarking that the results of
RobotCub will be fully open and consequently licensed following a
General Public (GP) license to the scientific community.
\end{quote}

\begin{enumerate}
\def\labelenumi{\arabic{enumi}.}
\setcounter{enumi}{15}
\item
  \ul{INTELTRAF}
\end{enumerate}

\begin{quote}
Title: Automatic Monitoring of Traffic and Detection of Accidents in
Highways
\end{quote}

Dates: 2002-12-01 to 2005-01-31 (successfully concluded)

Funding: ADI, POSI Med. 1.3 (22K€)

\begin{quote}
Institutions: \emph{Observit, Tecnologias de Visão por Computador Lda}

\emph{Instituto de Sistemas e Robótica, Instituto Superior Técnico,}

\emph{Aitek SRL} (IT)

\emph{Brisa, Auto-Estradas de Portugal.}
\end{quote}

Role: Deputy Coordinator at Participating Institution

\begin{quote}
Abstract: This project developed an Automated Traffic Surveillance
system with Computer Vision techniques, including real-time algorithms
for video sequence analysis of traffic scenes and the use of innovative
extended field of view camera systems (panoramic), for price reduction
and performance gain with respect to currently existing systems. We
address the applications of measuring traffic flow and detecting
abnormal events on highways and critical urban areas. Traffic flow
monitoring records statistical data on traffic distribution along time
(number of vehicles, average velocity, average wait time on queues,
accidents, transgressions to traffic rules, etc.).
\end{quote}

\begin{enumerate}
\def\labelenumi{\arabic{enumi}.}
\setcounter{enumi}{16}
\item
  \ul{CAVIAR}
\end{enumerate}

\begin{quote}
Title: Context Aware Vision using Image-based Active Recognition
\end{quote}

Dates: 2002-10-01 to 2005-09-30 (successfully concluded)

Funding: EU-FP5-IST-2001-37540 (1168K€/189K€)

\begin{quote}
Institutions: \emph{The University of Edimburgh,} (SC)

\emph{Instituto de Sistemas e Robótica, Instituto Superior Técnico}

\emph{Centre National de la Recherche Scientifique} (FR)

\emph{Institut National de Recherche en Informatique et en Automatique}
(FR)

\emph{Universite Joseph Fourier Grenoble 1} (FR)
\end{quote}

Role: Workpackage Coordinator at Participating Institution

\begin{quote}
Abstract: The main objective is to develop the theory of context-aware
visual recognition systems. We will implement the theory in a complete
closed-loop vision system, and apply it to two applications (city street
surveillance and customer behaviour analysis). To achieve these
objectives, we will develop new feature grouping, attention and
appearance-based recognition processes. This will also require
development of new techniques for acquiring, representing, and using
visual context and situation knowledge.
\end{quote}

\begin{enumerate}
\def\labelenumi{\arabic{enumi}.}
\setcounter{enumi}{17}
\item
  \ul{MIRROR}
\end{enumerate}

\begin{quote}
Title: Mirror Neurons for Recognition
\end{quote}

Dates: 2001-09-01 to 2004-08-31 (successfully concluded)

Funding: EU-FP5-IST-2000-28159 (720K€/160K€)

\begin{quote}
Institutions: \emph{Universita degli studi di Genova, Italy} (IT)

\emph{Universita degli studi di Ferrara} (IT)

\emph{Instituto de Sistemas e Robótica, Instituto Superior Técnico}

\emph{Uppsala Universitet} (SW)
\end{quote}

Role: Workpackage Coordinator at Participating Instituion

\begin{quote}
Abstract: The goals of MIRROR are (i) to realize an artificial system
that learns to communicate with humans by means of body gestures and
(ii) to study the mechanisms used by the brain to learn and represent
gestures. The biological base is the existence in primates pre-motor
cortex of a motor resonant system, called mirror neurons, activated both
during execution of goal directed actions and during observation of
similar actions performed by others. This unified representation may sub
serve the learning of goal directed actions during development and the
recognition of motor acts, when visually perceived.
\end{quote}

\begin{enumerate}
\def\labelenumi{\arabic{enumi}.}
\setcounter{enumi}{18}
\item
  \ul{RESCUE}
\end{enumerate}

\begin{quote}
Title: Cooperative Navigation for Rescue Robots
\end{quote}

Dates: 2001-09-01 to 2004-08-31 (successfully concluded)

Funding: FCT SRI/32546/99-00 (115K€)

\begin{quote}
Institutions: \emph{Instituto de Sistemas e Robótica, Instituto Superior
Técnico}
\end{quote}

Role: Team Member

\begin{quote}
Abstract: This project fosters research on multi-agent robotic systems
for search and rescue operations as its long-term goal. Currently, the
project is focused on obtaining new results on outdoors perception and
navigation, both for individual and cooperative robots.
\end{quote}

\begin{enumerate}
\def\labelenumi{\arabic{enumi}.}
\setcounter{enumi}{19}
\item
  \ul{NARVAL}
\end{enumerate}

\begin{quote}
Title: Navigation of Autonomous Robots Via Active Environmental
Perception
\end{quote}

Dates: 1996-03-01 to 2001-12-31 (successfully concluded)

Funding: EU-FP5-Esprit LTR Project -30185

\begin{quote}
Institutions: \emph{Instituto de Sistemas e Robótica, Instituto Superior
Técnico}

\emph{Université de Nice Sophia Antipolis, I3S-CNRS} (FR)

\emph{Universita degli studi di Genova} (IT)

\emph{Thomson Sintra ASM} (FR)
\end{quote}

Role: Responsible for Cross-platform Computational Architecture

\begin{quote}
Abstract: The goal of this project is to develop non-intruding and
reliable navigation systems giving the robot the ability to select
natural landmarks, and to navigate with respect to them, extending in
this way the autonomy range of autonomous robots operating in unknown
and unstructured environments. Reliability is achieved by continuously
controlling the uncertainty associated with knowledge of the environment
and of the robot's position and orientation. In the context of the
project, non-intrusiveness means that the robot must be able to operate
without special purpose landmarks being added to its environment.
Non-intrusive operation in unstructured environments precludes
navigation with respect to a set of handcrafted landmarks and requires
the robot's ability to infer its position from learned natural landmarks
of the environment using its perception system.
\end{quote}

\begin{enumerate}
\def\labelenumi{\arabic{enumi}.}
\setcounter{enumi}{20}
\item
  \ul{SIVA}
\end{enumerate}

\begin{quote}
Title: Integrated System for Active Surveillance
\end{quote}

Dates: 1998-01-01 to 2001-04-30 (successfully concluded)

Funding: FCT - PRAXIS 2/2.1/TPAR/2074/95 \hl{}

\begin{quote}
Institutions: \emph{Instituto de Sistemas e Robótica, Instituto Superior
Técnico}

\emph{Instituto de Sistemas e Robótica, Universidade de Coimbra}
\end{quote}

Role: Team Member

\begin{quote}
Abstract: With this project we aim to develop an autonomous system able
to perform surveillance tasks in a structured environment. The system
will be made up of several mobile platforms (agents). Each of them will
be equipped with an active vision system, an odometry system and
communication links enabling the platforms to communicate among
themselves. The full system (including all the platforms) will have a
high degree of autonomy. Each of the platforms will also be autonomous
in terms of the navigation and surveillance tasks. The system will be
designed to operate off-hours in structured environments such as
supermarkets, military installations, public buildings, power plants,
etc. The implication of its use during off-hours is that the primary
condition for the detection of an intruder will be the occurrence of
motion. Each one of the agents will perform the surveillance tasks based
on the active vision system.
\end{quote}

\begin{enumerate}
\def\labelenumi{\arabic{enumi}.}
\setcounter{enumi}{21}
\item
  \ul{INFANTE}
\end{enumerate}

\begin{quote}
Title: Development of Vehicles and Advanced Systems for the Execution of
Underwater Inspection Tasks
\end{quote}

Dates: 1997-01-01 to 2001-12-31 (successfully concluded)

Funding: FCT - PRAXIS XXI 3/3.1/TPAR/2042/95

\begin{quote}
Institutions: \emph{Instituto de Sistemas e Robótica, Instituto Superior
Técnico}

\emph{RINAVE}

\emph{Instituto Hidrográfico}

\emph{Cintal -- Universidade do Algarve}
\end{quote}

Role: Team Member

\begin{quote}
Abstract: The objectives of this project are the design and the
construction of an autonomous underwater vehicle (AUV), and the
development and integration of advanced systems for navigation, guidance
and control, acoustic communications, and mission management. These
systems will be tested in the lab. Once they are installed in the
vehicle, they will be also tested in pool and sea trials.
\end{quote}
